Leichte Sprache ist ein zentrales Instrument der barrierefreier Kommunikation und folgt klar definierten sprachlichen, inhaltlichen und formalen Regeln \cite{maas_leichte_2015}. 
Die automatisierte Erstellung von qualitativ hochwertigen Texten in Leichter Sprache stellt eine bedeutende Herausforderung dar. 
Sprachliche Reduktion, inhaltliche Vollständigkeit und Verständlichkeit müssen sorgfältig ausbalanciert werden, um den Bedürfnissen der Zielgruppe gerecht zu werden.

Im Kontext der beim Südwestrundfunk (SWR) eingesetzten Pipeline zur Erstellung von Texten in Leichter Sprache Plus (LS+) existieren bereits Module zur Faktenextraktion, KI-gestützen Übersetzung sowie zur regel- und LLM-basierten Validierung. 
Was jedoch fehlt, ist ein systematisches, automatisiertes Verfahren zur objektiven Bewertung der Textqualität sowie ein Mechanismus, der unzureichende Ergebisse erkennt und iterativ verbessert. 
Ohne einen solchen Ansatz ist die Qualität der generierten Texte nicht konstant und nur bedingt verlässlich nutzbar.